\documentclass[a4paper,12pt]{jlreq}
\usepackage{mylualatexstyle}

\title{Note}
\author{}
\date{}

\begin{document}
\maketitle
\begin{equation*}
    \hat{\mathcal{H}}_k = e^{-ik\cdot \hat{r}} \mathcal{H} e^{ik\cdot \hat{r}}
\end{equation*}
\begin{equation*}
    \hat{\mathcal{H}}_k \ket{u_n (k)} = E_n(k) \ket{u_n(k)}
\end{equation*}
\begin{align*}
    Q^{ij}_{n}(k)   & = \bra{\partial_i u_n(k)}\qty(1-\ket{u_n(k)}\bra{u_n(k)})\ket{\partial_j u_n(k)}\\
                & = \sum_{m \neq n} \frac{\bra{u_n(k)}\partial_i H_k \ket{u_m(k)}\bra{u_m(k)}\partial_j H_k \ket{u_n(k)}}{(E_m-E_n)^2}\\
                & = g^{ij}_{n}(k) - \frac{i}{2}\Omega^{ij}_{n}(k)
\end{align*}
\begin{align*}
    \dd s^2    & \coloneq 1-|\braket{u_n(k)}{u_n(k+\dd{k})}|^2\\
                & = g_{ij}(k)\dd k_i \dd k_j
\end{align*}
\begin{equation*}
    \gamma = \int_S \dd \bm{S} \cdot \bm{\Omega}_n(k) 
\end{equation*}

\newpage
\begin{equation*}
    a^{i}_{mn}(k) = i \braket{u_m(k)}{\partial_i u_n(k)}
\end{equation*}
\begin{equation*}
    \Omega^{ij}_n(k) = \partial_i a^{j}_{nn}(k) - \partial_j a^{i}_{nn}(k)
\end{equation*}
P Symmetry : 適切な Gauge をとると
\begin{equation*}
    P\ket{u_n(k)} = \ket{u_n(-k)}
\end{equation*}
\begin{equation*}
    a^{i}_{mn}(-k) = - a^{i}_{mn}(k) 
\end{equation*}
\begin{equation*}
    \Omega^{ij}_n(-k) =  \Omega^{ij}_n(k)
\end{equation*}

T Symmetry : 
\begin{equation*}
    T\ket{u_n(k)} = \ket{u_n(-k)}
\end{equation*}
\begin{equation*}
    a^{i}_{mn}(-k) = a^{i}_{nm}(k)
\end{equation*}
\begin{equation*}
    \Omega^{ij}_n(-k) =  -\Omega^{ij}_n(k)
\end{equation*}

PT Symmetry :
\begin{equation*}
    PT\ket{u_n(k)} = \ket{u_n(k)}
\end{equation*}
\begin{equation*}
    a^{i}_{mn}(k) = -a^{i}_{nm}(k)
\end{equation*}
\begin{equation*}
    \Omega^{ij}_n(k) =  0
\end{equation*}

\begin{equation*}
    (\braket{T\phi}{T\psi} = \braket{\psi}{\phi})
\end{equation*}

\newpage
\large{対称性とエネルギー縮退}

ハミルトニアン$\hat{H}$が群$G$の対称性を持つ, $g$に対応する表現を$\hat{C}(g)$とする.
$V = \mathrm{Span}\qty{\ket{\psi_1}, \cdots ,\ket{\psi_n}}$が群$G$の既約表現であるとは,
\begin{equation*}
    \forall g \in G,\ \hat{C}(g) V \subset V
\end{equation*}
を満たし, 以下のような部分空間$W \subset V$が存在しないことである。
\begin{equation*}
    \forall g \in G,\ \hat{C}(g) W \subset W
\end{equation*}

軌道混成:Schurの補題より, $V$上において$\hat{H}=E\hat{I}$となるため, エネルギーは空間$V$で$n$重縮退する.

\newpage
\large{回転対称性}

波数$\vec{k}\in \mathrm{BZ}$について, 小群$G_k \subset G$は以下を満たす
\begin{equation*}
    \forall g \in G_k,\ g(\vec{k}) \equiv \vec{k}
\end{equation*}
このとき$\mathcal{H}_k$と$G_k$は可換なので
$V = \mathrm{Span}\qty{\ket{u_1(\vec{k})}, \cdots ,\ket{u_n(\vec{k})}}$が$G_k$の既約表現なら
$\vec{k}$点で$n$重縮退する
\end{document}